\chapter{Verhaltensmodellierung mit Activity Diagrams}

\begin{lernziele}
Nach diesem Kapitel können Sie Abläufe des Roboters als Aktivitätsdiagramm beschreiben.
\end{lernziele}

\section{Was beschreibt ein Aktivitätsdiagramm?}
Ein Aktivitätsdiagramm beschreibt einen Ablauf aus Aktionen, Entscheidungen und Verzweigungen. Es eignet sich gut für Prozesse wie \enquote{Navigationsziel anfahren}.

\section{Beispielablauf Navigation}
\begin{enumerate}
\item System starten
\item Sensoren initialisieren
\item Karte laden
\item Ziel empfangen
\item Pfad planen
\item Bewegung ausführen
\item Hindernisse prüfen
\item Ziel erreicht melden
\end{enumerate}

\section{Entscheidungen}
Entscheidungsknoten modellieren Alternativen. Beispiel: Wenn ein Hindernis erkannt wird, plant der Roboter neu oder stoppt.

\section{Parallelität}
Ein Roboter führt oft mehrere Aktivitäten parallel aus. Während er fährt, überwacht er Sensoren und Batterie. Aktivitätsdiagramme können solche parallelen Abläufe darstellen, sollten am Anfang aber nicht überladen werden.

\begin{praxisbox}
Der Ablauf \enquote{Zur Ladestation fahren} enthält eine Entscheidung: Ist der Akkustand niedrig? Nur dann wird der Ladevorgang gestartet.
\end{praxisbox}

\begin{uebungbox}
Modellieren Sie den Ablauf \enquote{Ziel anfahren}. Ergänzen Sie mindestens eine Entscheidung für Hinderniserkennung und eine für niedrigen Akkustand.
\end{uebungbox}
